%%% SECTION 11: NATIONAL MCP SERVER INFRASTRUCTURE %%%

Physical AI National MCP Servers \cite{national-mcp-paper} provides the
national-scale infrastructure for connecting Physical AI oncology trial systems
across sites using the Model Context Protocol (MCP) \cite{mcp-protocol}. The
five-server architecture addresses the fragmentation problem that prevents
current clinical trial systems from interoperating, establishing standardized
communication pathways for Physical AI operations at scale. The infrastructure
is implemented in the National MCP Servers Repository \cite{mcp-repo} with 23
tools, conformance levels, safety modules, and deployment configurations.

\subsection{The Fragmentation Problem}
\label{subsec:mcp-fragmentation}

Current clinical trial systems operate in silos with incompatible data formats,
disconnected communication channels, and no standardized protocol for real-time
inter-system coordination. Electronic data capture (EDC) systems, laboratory
information management systems (LIMS), clinical trial management systems (CTMS),
and regulatory submission platforms each use proprietary data formats and
communication protocols. ClinicalTrials.gov \cite{clinicaltrials-gov} provides
trial registration but not real-time operational coordination. FDA clinical
trials guidance \cite{fda-clinical-trials-guidance} addresses individual system
requirements but not inter-system communication.

Physical AI magnifies this fragmentation because robots, AI systems, clinical
databases, and regulatory reporting tools must all communicate seamlessly in
real time. A surgical robot generating telemetry data must share that data with
monitoring systems, clinical databases, regulatory reporting tools, and safety
interlock systems simultaneously. FHIR \cite{hl7-fhir} addresses healthcare
data exchange and DICOM \cite{dicom} addresses medical imaging, but neither
provides a complete solution for the real-time, multi-modal data exchange
required by Physical AI systems.

\subsection{Model Context Protocol for Physical AI}
\label{subsec:mcp-protocol}

The Model Context Protocol \cite{mcp-protocol} provides a standardized
framework for AI systems to communicate with tools and data sources. Originally
designed for language model tool use, MCP is adapted for Physical AI oncology
trial operations by extending the protocol to handle robot control commands,
sensor data streams, safety interlock signals, and regulatory reporting
messages. The prior architecture that the five-server design builds upon
includes individual MCP server implementations for robot control, data
management, and safety monitoring that were developed in earlier versions of
the Physical AI platform \cite{main-repo}.

The adaptation of MCP for Physical AI establishes several key capabilities:
bidirectional communication between AI planning systems and physical robots,
standardized data schemas for robot telemetry and sensor readings, real-time
safety signal propagation across all connected systems, and audit trail
generation compliant with 21 CFR Part 11 electronic records requirements.

\subsection{Five-Server Architecture}
\label{subsec:mcp-architecture}

The national MCP infrastructure comprises five specialized servers, each
responsible for a distinct aspect of Physical AI trial operations.

\begin{table}[H]
\centering
\caption{Five-Server MCP Architecture for Physical AI Oncology Trials}
\label{tab:mcp-servers}
\small
\begin{tabularx}{\textwidth}{c l X}
\toprule
\textbf{No.} & \textbf{Server} & \textbf{Responsibilities} \\
\midrule
1 & Governance Server & Trial protocol management, regulatory compliance tracking, role-based access control, audit trail management \\
2 & Data Flow Server & Real-time data routing, format transformation, data validation, storage coordination across FHIR and DICOM \\
3 & Safety Server & Emergency stop propagation, safety interlock management, anomaly detection, incident reporting \\
4 & Interoperability Server & Cross-platform translation, device registration, capability negotiation, standards conformance \\
5 & Analytics Server & Real-time dashboards, trend analysis, outcome reporting, predictive analytics for trial management \\
\bottomrule
\end{tabularx}
\end{table}

The Governance Server manages trial protocol configurations, ensuring that all
Physical AI operations comply with the approved protocol and the three adapted
regulatory standards \cite{ich-adapt, cfr50-adapt, cfr312-adapt}. The Data
Flow Server handles the high-volume, multi-modal data streams generated by
Physical AI systems, routing data to appropriate destinations based on type,
urgency, and regulatory requirements. The Safety Server provides the most
time-critical function: propagating emergency stop signals across all connected
Physical AI systems within milliseconds, ensuring that a safety event detected
by any system triggers appropriate responses across the entire infrastructure.

The Interoperability Server addresses the challenge of connecting diverse
Physical AI systems from different manufacturers with different communication
protocols. It provides translation services, device registration, and
capability negotiation, ensuring that a Franka Panda cobot can exchange data
with a da Vinci surgical robot despite different native protocols. The
Analytics Server aggregates data from all other servers to provide real-time
operational dashboards, trend analysis, and predictive analytics that support
trial management decisions.

\subsection{Tool Inventories and Conformance Levels}
\label{subsec:mcp-tools}

The five servers collectively provide 23 tools organized into functional
categories. Each tool has a defined conformance level specifying the minimum
requirements for implementations. Conformance levels range from Level 1 (basic
data exchange) through Level 3 (full bidirectional integration with audit trail
and safety interlock support). The MCP repository \cite{mcp-repo} provides
reference implementations for all 23 tools at Level 2 or above.

Tool categories include robot management tools (registration, status query,
command execution, calibration verification), data management tools (data
ingestion, transformation, validation, archival), safety tools (emergency stop,
interlock query, anomaly reporting, incident documentation), compliance tools
(protocol verification, consent tracking, reporting generation, audit trail
query), and analytics tools (dashboard generation, trend computation, outcome
aggregation, performance scoring).

\subsection{Safety Modules and Emergency Systems}
\label{subsec:mcp-safety}

The safety architecture embedded in the MCP server infrastructure provides
multiple layers of protection. The emergency stop system propagates stop signals
from any source (hardware button, software trigger, automated anomaly
detection) to all connected Physical AI systems within a specified latency
threshold. The automatic failover system ensures that safety functions continue
even if individual server components fail, using redundant safety modules that
operate independently of the primary server infrastructure.

Safety interlock protocols ensure that Physical AI systems cannot perform
restricted operations without proper authorization. For example, a surgical
robot cannot initiate a procedure until the Governance Server confirms that
the patient has provided informed consent, the Safety Server confirms that all
pre-procedure safety checks have passed, and the Data Flow Server confirms that
all monitoring systems are recording. This multi-server interlock provides
defense-in-depth that exceeds the safety capabilities of any single system.

\subsection{Integration and Interoperability}
\label{subsec:mcp-integration}

Integration adapters connect the MCP infrastructure to existing healthcare
systems. FHIR \cite{hl7-fhir} adapters enable data exchange with electronic
health record systems. DICOM \cite{dicom} adapters handle medical imaging data
from Physical AI imaging robots. HL7 v2 adapters support legacy hospital
information systems. RESTful API adapters enable integration with modern
web-based clinical trial management platforms.

The interoperability testbed validates that MCP servers can communicate
reliably with these external systems under operational conditions. Test
coverage includes data format conversion accuracy, message delivery latency,
error handling and recovery, and security boundary enforcement. SDK and
tooling packages enable third-party developers to build MCP-compatible
Physical AI applications.

\subsection{National Deployment Topology}
\label{subsec:mcp-deployment}

The national deployment architecture uses a hub-and-spoke model. Regional hub
servers aggregate data from multiple trial sites within a geographic region,
providing regional analytics and safety coordination. Spoke servers at
individual trial sites handle local Physical AI operations. A national
coordination layer connects all regional hubs, providing nationwide analytics,
cross-regional safety signal aggregation, and centralized regulatory reporting.

Data routing ensures that patient-identified data remains within authorized
boundaries (site-level and regional as required by HIPAA \cite{hipaa} and state
privacy laws), while de-identified aggregate data flows to the national level
for trend analysis and quality improvement. The federated learning framework
\cite{fl-paper} interfaces with the MCP infrastructure at the regional hub
level, enabling model improvement without centralizing patient data.

\subsection{Governance and Compliance}
\label{subsec:mcp-governance}

The governance framework for national MCP server operations addresses server
management authority, update deployment procedures, compliance monitoring
protocols, and dispute resolution mechanisms. Server management is distributed
across three levels: site-level administration for local configurations,
regional administration for hub coordination, and national administration for
infrastructure-wide updates and standards enforcement.

Compliance monitoring uses the Governance Server's audit trail capabilities to
track all Physical AI operations against regulatory requirements. Automated
compliance checks flag potential violations in real time, enabling corrective
action before regulatory deadlines. HIPAA \cite{hipaa} compliance is enforced
through access controls, encryption, and audit trails managed by the
Governance and Data Flow servers.

\subsection{CI/CD and Operational Excellence}
\label{subsec:mcp-cicd}

The continuous integration and deployment pipeline for MCP servers ensures
reliable updates with minimal operational disruption. Automated testing covers
unit tests, integration tests, and end-to-end scenario tests for all 23 tools
and five servers. Deployment procedures support rolling updates that maintain
service availability during version transitions. Rollback capabilities enable
rapid reversion to previous versions if issues are detected post-deployment.
The main repository \cite{main-repo} and MCP repository \cite{mcp-repo}
provide the CI/CD configurations and test suites.

\subsection{Patient Safety and Data Flow}
\label{subsec:mcp-patient-safety}

MCP servers ensure patient safety through real-time monitoring, data
validation, and automated safety checks. The complete data flow from robot
procedures through MCP servers to clinical databases and regulatory reporting
systems maintains data integrity at every step. Each data element is validated
against expected ranges, flagged for anomalies, and logged with timestamps and
source identification.

The integration between MCP servers and the federated learning framework
\cite{fl-paper} enables multi-site model improvement while preserving patient
privacy. Model updates generated through federated learning are distributed
through MCP channels, with the Safety Server verifying that updates meet safety
requirements before deployment to Physical AI systems. This creates a closed
loop from data collection through model improvement to deployment, all managed
through standardized MCP infrastructure that supports federated learning in
healthcare \cite{federated-learning-healthcare}.

\subsection{Complete Tool Inventory by Server}
\label{subsec:mcp-tool-inventory}

Table~\ref{tab:mcp-all-tools} provides the complete inventory of all 23 MCP
tools organized by their parent server, with each tool's name and functional
description. This inventory serves as the authoritative reference for
implementers deploying the five-server architecture at trial sites.

\begin{longtable}{c l l p{6.5cm}}
\caption{Complete Inventory of 23 MCP Tools Organized by Server}
\label{tab:mcp-all-tools} \\
\toprule
\textbf{No.} & \textbf{Server} & \textbf{Tool Name} & \textbf{Description} \\
\midrule
\endfirsthead
\toprule
\textbf{No.} & \textbf{Server} & \textbf{Tool Name} & \textbf{Description} \\
\midrule
\endhead
\midrule
\multicolumn{4}{r}{\textit{Continued on next page}} \\
\bottomrule
\endfoot
\bottomrule
\endlastfoot
1 & Governance & Protocol Manager & Manages trial protocol configurations, version control, and distribution to connected systems \\
2 & Governance & Compliance Tracker & Monitors regulatory compliance status across all three adapted standards in real time \\
3 & Governance & Access Controller & Enforces role-based access control policies aligned with HIPAA requirements and 21 CFR Part 11 \\
4 & Governance & Audit Logger & Records all system operations with timestamps, user identities, and action details for regulatory audit trails \\
5 & Governance & Consent Verifier & Validates patient informed consent status before authorizing Physical AI procedures \\
\midrule
6 & Data Flow & Data Ingester & Receives and normalizes incoming data streams from Physical AI systems, sensors, and clinical databases \\
7 & Data Flow & Format Transformer & Converts data between FHIR, DICOM, HL7 v2, and internal Physical AI data schemas \\
8 & Data Flow & Data Validator & Applies range checks, consistency validation, and completeness verification to all incoming data elements \\
9 & Data Flow & Storage Coordinator & Routes validated data to appropriate storage systems based on data type, retention policy, and access requirements \\
10 & Data Flow & Stream Router & Manages real-time data routing based on urgency level, data type, and destination system requirements \\
\midrule
11 & Safety & Emergency Stop & Propagates emergency stop signals from any source to all connected Physical AI systems within latency threshold \\
12 & Safety & Interlock Manager & Manages safety interlock states ensuring that Physical AI systems cannot perform restricted operations without authorization \\
13 & Safety & Anomaly Detector & Monitors Physical AI telemetry streams for anomalous patterns indicating potential safety concerns \\
14 & Safety & Incident Reporter & Generates structured incident reports for safety events compliant with adapted ICH E6(R3) adverse event requirements \\
\midrule
15 & Interopertic & Device Registrar & Registers Physical AI devices with capability profiles, communication protocols, and safety certifications \\
16 & Interopertic & Protocol Translator & Translates commands and data between different Physical AI device communication protocols \\
17 & Interopertic & Capability Negotiator & Negotiates operational parameters between systems with different capabilities and performance characteristics \\
18 & Interopertic & Standards Enforcer & Validates that all cross-platform communications conform to MCP protocol specifications and conformance levels \\
19 & Interopertic & Integration Adapter & Provides connectivity adapters for FHIR, DICOM, HL7 v2, and RESTful API integration with external systems \\
\midrule
20 & Analytics & Dashboard Generator & Creates real-time operational dashboards displaying trial status, system health, and safety metrics \\
21 & Analytics & Trend Analyzer & Computes trend analyses across patient outcomes, system performance, and adverse event patterns \\
22 & Analytics & Outcome Aggregator & Aggregates treatment outcome data across patients and procedures for regulatory reporting and quality improvement \\
23 & Analytics & Performance Scorer & Calculates performance scores for Physical AI systems based on operational metrics and clinical outcomes \\
\end{longtable}

Each tool operates at a defined conformance level that specifies the minimum
requirements for compliant implementations. Level 1 conformance requires basic
data exchange capability with standard message formats. Level 2 conformance
adds bidirectional communication, error handling, retry logic, and basic audit
trail support. Level 3 conformance, the highest level, requires full
bidirectional integration with comprehensive audit trail generation, safety
interlock participation, and real-time monitoring integration. The reference
implementations in the MCP repository \cite{mcp-repo} achieve Level 2 or above
for all 23 tools, with the Safety Server tools (Emergency Stop, Interlock
Manager, Anomaly Detector, Incident Reporter) all implemented at Level 3 given
their critical role in patient safety.

The tool inventory reflects a deliberate design decision to distribute
functionality across servers rather than centralizing it. This distribution
provides fault isolation: a failure in the Analytics Server does not affect
safety operations managed by the Safety Server. It also supports independent
scaling, allowing sites with higher data volumes to scale Data Flow Server
resources without modifying other server components. The Governance Server tools
(Protocol Manager, Compliance Tracker, Access Controller, Audit Logger, Consent
Verifier) interact with every other server, making the Governance Server the
coordination hub of the architecture. The five tools on the Governance Server
collectively ensure that no Physical AI operation proceeds without proper
authorization, consent verification, and audit documentation.

The Interoperability Server carries five tools because cross-platform
translation is the most complex operational challenge in a multi-vendor Physical
AI environment. The Device Registrar maintains a registry of all connected
Physical AI systems with their capability profiles, enabling the Protocol
Translator to select appropriate translation strategies for each device pair.
The Capability Negotiator resolves conflicts when two systems have different
performance characteristics, such as when a high-frequency telemetry source
connects to a lower-bandwidth monitoring system. The Standards Enforcer
validates all cross-platform communications against MCP protocol specifications,
rejecting non-conformant messages before they reach destination systems. The
Integration Adapter provides the bridge to external healthcare systems,
supporting FHIR \cite{hl7-fhir} for electronic health records, DICOM
\cite{dicom} for medical imaging, HL7 v2 for legacy hospital systems, and
RESTful APIs for modern web-based platforms.

\subsection{National Deployment Phases}
\label{subsec:mcp-deployment-phases}

The national deployment of MCP server infrastructure proceeds through four
phases, each expanding geographic coverage and operational capability. The
phased approach ensures that lessons learned at each stage inform subsequent
deployments, reducing risk and improving deployment efficiency.

\begin{longtable}{l p{2.2cm} p{2.5cm} p{2.2cm} p{4.0cm}}
\caption{National MCP Server Deployment Phases}
\label{tab:mcp-deployment-phases} \\
\toprule
\textbf{Phase} & \textbf{Timeline} & \textbf{Geographic Scope} & \textbf{Sites} & \textbf{Infrastructure Configuration} \\
\midrule
\endfirsthead
\toprule
\textbf{Phase} & \textbf{Timeline} & \textbf{Geographic Scope} & \textbf{Sites} & \textbf{Infrastructure Configuration} \\
\midrule
\endhead
\bottomrule
\endfoot
\bottomrule
\endlastfoot
Phase A: Prototype & Months 1-6 & California only & 1-2 sites & Single hub server, local spoke servers, direct connections \\
\midrule
Phase B: Regional & Months 7-18 & West Coast states & 5-10 sites & Regional hub server for West Coast, spoke servers at each site, federated learning node activation \\
\midrule
Phase C: Multi-Regional & Months 19-36 & Four U.S. regions (West, South, Midwest, Northeast) & 20-50 sites & Four regional hub servers, national coordination layer, cross-regional safety signal aggregation \\
\midrule
Phase D: Nationwide & Months 37-60 & All 50 states plus territories & 100+ sites & Full hub-and-spoke topology, redundant national coordination, complete federated learning network \\
\end{longtable}

Phase A (Prototype) deploys the five-server architecture at the initial
California trial site or sites described in
Section~\ref{sec:implementation}. This phase validates the MCP infrastructure
under real operational conditions, comparing actual performance against the
metrics established during the 24-hour simulation \cite{site-docs}. A single
hub server manages all local spoke servers, and direct connections between
servers minimize network complexity during initial validation. The prototype
phase produces operational metrics including message latency, data throughput,
safety signal propagation time, and system availability that inform subsequent
deployment phases.

Phase B (Regional) expands MCP infrastructure to additional West Coast sites,
deploying a regional hub server that aggregates data from all sites within the
region. This phase introduces the hub-and-spoke topology at regional scale,
validating the architecture's ability to coordinate multiple sites through a
shared hub. Federated learning nodes activate at each site, enabling cross-site
model improvement through the framework described in
Section~\ref{sec:federated-learning}. Network requirements increase
significantly in Phase B as inter-site communication adds bandwidth and latency
constraints that were absent in the single-site prototype.

Phase C (Multi-Regional) establishes four regional hub servers covering the
major U.S. geographic regions. The national coordination layer connects all
regional hubs, enabling nationwide analytics, cross-regional safety signal
aggregation, and centralized regulatory reporting. This phase introduces the
most complex operational challenge: managing data routing across multiple
regional boundaries while maintaining compliance with varying state privacy
laws. The Data Flow Server's Storage Coordinator tool becomes critical in
Phase C, enforcing data residency requirements that differ across states. HIPAA
\cite{hipaa} provides a federal baseline, but states such as California (CMIA,
CCPA) impose additional requirements that the MCP infrastructure must respect
at the routing level.

Phase D (Nationwide) achieves full geographic coverage with MCP infrastructure
deployed in all 50 states plus U.S. territories. The infrastructure supports
100 or more connected trial sites with redundant national coordination ensuring
continuity even if individual regional hubs experience outages. The complete
federated learning network enables model improvement across the full national
dataset while preserving patient privacy at every level. Phase D represents the
steady-state operational configuration described in the hub-and-spoke model,
with ongoing expansion accommodating new trial sites as they achieve PSL
compliance and activate Physical AI operations.

Each deployment phase includes validation gates that must be passed before
proceeding to the next phase. Phase A requires demonstration of all 23 tools
operating at Level 2 conformance or above under operational conditions. Phase B
requires successful cross-site data exchange with latency within specified
thresholds and federated learning round completion across all participating
sites. Phase C requires cross-regional safety signal propagation within latency
requirements and compliance with all applicable state privacy laws. Phase D
requires sustained operation at nationwide scale with 99.5 percent minimum
uptime across all regional hubs and the national coordination layer.

The deployment timeline is conservative by design, allocating five years for
full nationwide coverage. This timeline accommodates the legislative processes
required in each state (analogous to California's SB 1042, AB 2847, and SB 892
from the site documentation \cite{site-docs}), the infrastructure construction
and validation activities at each site, and the workforce training requirements
described in Section~\ref{sec:implementation}. Sites in states with existing
legislative frameworks for AI in healthcare may proceed through earlier phases
more quickly, while sites in states requiring new legislation may follow the
later phases of the timeline.
